% =====================================================================
% Plantilla del artículo del proyecto final
% Maestría en Inteligencia Artificial — Agentes, Skills y MCP
%
% Extensión: 6 a 8 páginas en este formato, SIN contar referencias ni apéndices.
% Compilar:  pdflatex paper && bibtex paper && pdflatex paper && pdflatex paper
% También compila sin cambios en Overleaf (subir paper.tex y referencias.bib).
%
% Cada sección trae un comentario con lo que se evalúa. Bórrenlos antes de entregar.
% =====================================================================
\documentclass[10pt,twocolumn,a4paper]{article}

\usepackage[utf8]{inputenc}
\usepackage[T1]{fontenc}
\usepackage[spanish]{babel}
\usepackage[margin=2cm]{geometry}
\usepackage{booktabs}
\usepackage{graphicx}
\usepackage{amsmath}
\usepackage{xcolor}
\usepackage[hidelinks]{hyperref}
\usepackage{listings}
\usepackage{caption}

\lstset{basicstyle=\ttfamily\footnotesize, breaklines=true, frame=single, columns=fullflexible}
\newcommand{\todo}[1]{\textcolor{red}{[PENDIENTE: #1]}}

\title{\todo{Título: qué se hizo y qué se encontró, no el nombre del sistema}}
\author{
  José Andrés Lizano-Gallegos \\
  \small Maestría en Inteligencia Artificial\\
  \small \texttt{jlizanog@ucenfotec.ac.cr}
}
\date{}

\begin{document}
\maketitle

% ---------------------------------------------------------------------
\begin{abstract}
% 150 a 200 palabras. Cuatro movimientos, una o dos frases cada uno:
% (1) problema y por qué importa, (2) qué proponen, (3) cómo lo evaluaron,
% (4) resultado principal CON UNA CIFRA. Sin citas, sin siglas sin definir.
\todo{Resumen en español.}
\end{abstract}

\noindent\textbf{Palabras clave:} \todo{4 a 6 palabras clave}

\noindent\textbf{Abstract.} \todo{Traducción al inglés del resumen.}

% ---------------------------------------------------------------------
\section{Introducción}
% Qué se evalúa: el problema está acotado a un usuario y un contexto concretos;
% queda claro por qué un agente (y no un flujo fijo) es la aproximación adecuada;
% las preguntas de investigación son explícitas y medibles; las contribuciones
% están enumeradas y cada una se sostiene en una sección del artículo.
\todo{Contexto y problema.}

\todo{Por qué un agente: qué parte del proceso no puede decidirse de antemano.}

Este trabajo responde las siguientes preguntas de investigación:
\begin{itemize}
  \item \textbf{PI1:} \todo{p. ej., ¿en qué medida una Skill procedimental mejora la exactitud frente a instrucciones genéricas?}
  \item \textbf{PI2:} \todo{p. ej., ¿qué costo en tokens y latencia introduce?}
  \item \textbf{PI3:} \todo{p. ej., ¿qué proporción de inyecciones indirectas logra alterar una acción?}
\end{itemize}

Las contribuciones son: \todo{(i) ..., (ii) ..., (iii) ...}

% ---------------------------------------------------------------------
\section{Trabajo relacionado}
% Mínimo 8 referencias; al menos 5 revisadas por pares (congreso o revista).
% No es un catálogo: agrupen por tema y digan en qué se parece y en qué se
% distingue su trabajo de cada grupo. TODA referencia debe haber sido abierta
% y leída por el equipo. Una cita inexistente o inventada anula esta sección.
\subsection{Agentes basados en modelos de lenguaje}
\todo{Razonamiento y acción intercalados \cite{yao2023react}, uso aprendido de herramientas \cite{schick2023toolformer}, autocorrección \cite{shinn2023reflexion}.}

\subsection{Protocolos de herramientas y contexto}
\todo{MCP \cite{mcp_spec} y alternativas; patrones de diseño de agentes \cite{anthropic2024effective}.}

\subsection{Evaluación y seguridad de agentes}
\todo{Bancos de pruebas de agentes \cite{liu2024agentbench,jimenez2024swebench}; inyección de prompts indirecta \cite{greshake2023injection}.}

% ---------------------------------------------------------------------
\section{Sistema propuesto}
% Qué se evalúa: la arquitectura se entiende con una figura; cada decisión de
% diseño está justificada (no solo descrita); la política de permisos y los
% límites del bucle están explícitos.
\subsection{Arquitectura}
\todo{Figura \ref{fig:arquitectura} con host, clientes, servidores MCP, Skill y política.}

\begin{figure}[t]
  \centering
  \fbox{\parbox{0.9\columnwidth}{\centering\vspace{2cm}\todo{diagrama de arquitectura}\vspace{2cm}}}
  \caption{Arquitectura del agente. \todo{Describir qué muestra la figura.}}
  \label{fig:arquitectura}
\end{figure}

\subsection{Servidores MCP y herramientas}
\todo{Tabla con cada herramienta: servidor, lectura o escritura, validaciones del lado del servidor.}

\subsection{Skill y procedimiento}
\todo{Qué conocimiento procedimental codifica la Skill y qué decide el código.}

\subsection{Política de permisos y control del bucle}
\todo{Herramientas permitidas, con confirmación, denegadas y ocultas; máximo de pasos; presupuesto.}

% ---------------------------------------------------------------------
\section{Metodología de evaluación}
% Qué se evalúa: el conjunto de evaluación es propio, tiene al menos 40 casos
% con al menos 30% difíciles, y fue etiquetado por dos personas con acuerdo
% reportado. Hay una línea base real. Cada configuración se ejecuta al menos
% 3 veces. Las métricas se definen ANTES de ver los resultados.
\subsection{Conjunto de evaluación}
\todo{Origen de los datos, cómo se anonimizaron, cuántos casos y de qué dificultad.}

\todo{Etiquetado por dos anotadores; acuerdo con kappa de Cohen \cite{cohen1960kappa}: $\kappa = $ valor.}

\subsection{Configuraciones comparadas}
\todo{Línea base vs. propuesta. Qué variable cambia y cuáles se mantienen fijas.}

\subsection{Métricas}
\begin{itemize}
  \item \textbf{Exactitud} o tasa de tarea completada: \todo{definición exacta}.
  \item \textbf{Costo:} tokens de entrada y salida por ejecución.
  \item \textbf{Latencia:} segundos por ejecución.
  \item \textbf{Eficiencia:} pasos y llamadas a herramientas.
  \item \textbf{Robustez:} proporción de ataques que alteran una acción o una cifra.
\end{itemize}

% ---------------------------------------------------------------------
\section{Resultados}
% Qué se evalúa: se reporta media y dispersión (no un solo número); las tablas
% responden a las preguntas de investigación en orden; no se interpreta aquí,
% se reporta. Toda cifra debe poder reproducirse desde las trazas entregadas.
\begin{table}[t]
  \centering
  \caption{Resultados por configuración (media $\pm$ desviación, $n=$\todo{}).}
  \label{tab:resultados}
  \small
  \begin{tabular}{lccc}
    \toprule
    Métrica & Línea base & Propuesta & $\Delta$ \\
    \midrule
    Exactitud (\%)        & \todo{} & \todo{} & \todo{} \\
    Tokens por ejecución  & \todo{} & \todo{} & \todo{} \\
    Latencia (s)          & \todo{} & \todo{} & \todo{} \\
    Pasos                 & \todo{} & \todo{} & \todo{} \\
    Ataques exitosos (\%) & \todo{} & \todo{} & \todo{} \\
    \bottomrule
  \end{tabular}
\end{table}

\subsection{PI1}
\todo{Resultado con referencia a la Tabla \ref{tab:resultados}.}
\subsection{PI2}
\todo{}
\subsection{PI3}
\todo{}

\subsection{Análisis de errores}
% Obligatorio. Clasifiquen los fallos (no los cuenten solamente) y muestren
% dos o tres ejemplos concretos tomados de las trazas.
\todo{Tipología de errores con ejemplos.}

% ---------------------------------------------------------------------
\section{Discusión}
% Qué significan los resultados, qué sorprendió, y qué recomendaría el equipo
% a alguien que construya un agente similar. Aquí sí se interpreta.
\todo{}

\section{Limitaciones y amenazas a la validez}
% Obligatorio y honesto: tamaño del conjunto, sesgo del etiquetado, dependencia
% del modelo y su versión, validez externa, qué no se probó.
\todo{}

\section{Consideraciones éticas}
% Datos personales y su anonimización, consentimiento, riesgos del sistema si
% se despliega, y quién es responsable de sus acciones.
\todo{}

\section{Conclusiones}
\todo{Respuesta directa a cada pregunta de investigación y trabajo futuro.}

% ---------------------------------------------------------------------
\section*{Declaración de uso de IA}
% Obligatoria. Qué herramientas de IA usaron para escribir, programar o analizar,
% para qué, y cómo verificaron lo que produjeron. Usar IA está permitido;
% no declararlo, no.
\todo{}

\section*{Disponibilidad del código y los datos}
\todo{Enlace al repositorio y commit exacto usado para los resultados.}

\bibliographystyle{plain}
\bibliography{referencias}

\appendix
\section{Reproducibilidad}
% Comandos exactos para regenerar la Tabla 1 desde cero; versión del modelo,
% fecha de ejecución, semillas, y dónde están las trazas.
\begin{lstlisting}
# Ejemplo: regenerar la tabla de resultados
python agente.py --config config_base.json      # x3
python agente.py --config config_propuesta.json # x3
python leer_traza.py trazas/*.jsonl --csv metricas.csv
\end{lstlisting}

\end{document}
